%% Chapter 4 — Prototyping / System Design
\label{ch:prototyping}

This chapter presents the EMS Portal as it is currently built. It documents the
architecture and technology stack, the feature modules completed in Phases~1 and~2,
the authorization model, the database schema with its migrations, and the quality
assurance practices. As of June~2026, the foundational architecture and base
modules are in place, ready for the remaining feature phases.

\section{System Architecture}
The EMS Portal is organized as a single repository containing three related
packages: a frontend (user-facing web application), a backend (API server), and
a shared types package that keeps both in sync. It follows a three-tier
client-server design: the presentation tier (browser interface), the business
logic tier (API server), and the data tier (database). Figure~\ref{fig:arch}
shows the layout.

\begin{figure}[ht]
	\centering
	\begin{tikzpicture}[
		node distance=0.9cm and 1.4cm,
		box/.style={draw, rounded corners, align=center, font=\footnotesize,
			minimum width=3.2cm, minimum height=1.0cm},
		tier/.style={draw, dashed, rounded corners, inner sep=10pt},
		arr/.style={-{Stealth[length=2mm]}, thick},
		lbl/.style={font=\scriptsize\itshape}
	]
		\node[box, fill=blue!5]   (browser) {Browser\\(Single-Page App)};
		\node[box, fill=blue!5]   (spa) [below=of browser] {React 19 + Vite\\TanStack Router/Query\\Jotai, Tailwind};
		\node[box, fill=green!6]  (api) [below=of spa] {NestJS API (Fastify)\\Controllers $\to$ Services\\$\to$ Repositories};
		\node[box, fill=orange!8] (db) [below=of api] {PostgreSQL\\(TypeORM, Data Mapper)};

		\node[box, fill=gray!8, minimum height=3.4cm, text width=2.6cm]
			(shared) [right=2.2cm of spa.east, yshift=-0.9cm]
			{\textbf{@ems-portal/}\\\textbf{types}\\[2pt]Zod schemas\\(shared DTOs)};

		\draw[arr, transform canvas={xshift=-3mm}] (browser) -- (spa);
		\draw[arr, transform canvas={xshift=3mm}] (spa) -- (browser);
		\draw[arr] (spa) -- node[lbl, right] {REST / HTTPS} (api);
		\draw[arr] (api) -- node[lbl, right] {SQL} (db);

		\draw[arr, dashed, gray] (shared.north) |- (spa.east);
		\draw[arr, dashed, gray] (shared.south) |- (api.east);
	\end{tikzpicture}
	\caption{Three-tier architecture of the EMS Portal. The shared types package
		is used by both the frontend and backend, keeping the network contract in
		sync at compile time.}
	\label{fig:arch}
\end{figure}

\subsection{Technology Stack}
The backend is a NestJS REST API running on Fastify (a fast HTTP server). It
persists data to PostgreSQL through TypeORM using the \textbf{Data Mapper}
pattern, a design where entities (database models) are kept separate from the
queries that fetch them, making the code modular and testable. Requests are
validated against Zod schemas (formal, machine-readable definitions of allowed
data shapes), and the API exposes auto-generated interactive documentation via
Swagger.

The frontend is a React~19 single-page application built with Vite (a build
tool). It uses TanStack Router for file-based routing (routes are defined as
files in a directory, not programmatically), TanStack Query to manage server
state as a cache layer (keeping the server's data in sync without manual
synchronization), Jotai for client-only state (UI toggles, form
drafts), and Tailwind CSS with RadixUI for styling and pre-built UI
components.

A key design principle is the shared types package (\texttt{@ems-portal/types}),
which defines the data structure for every piece of information exchanged between
the frontend and backend using Zod schemas. Both the frontend and backend
reference these shared definitions, so the compiler catches any inconsistencies
before the code runs, it acts as a critical safeguard for type safety across the network
boundary. Backend entities \texttt{implement} the inferred types, meaning if the
shared schema changes, any code that no longer matches gets a compiler error,
preventing silent bugs.

\subsection{Backend Module Pattern}
The backend is organized into feature modules, each handling a distinct piece of
functionality (authentication, user management, provider profiles, etc.). Each
module follows the same internal structure: a \textbf{controller} layer (handles
incoming HTTP requests and responses), a \textbf{service} layer (contains the
business logic), an injectable \textbf{repository} layer (encapsulates custom
database queries), and one or more \textbf{entities} extending a common
\texttt{BaseEntity} that automatically tracks record creation and modification
times. This consistent, layered structure makes the codebase predictable,
testable, and maintainable.

All requests are validated against the shared Zod schemas via a global
validation pipe before reaching the service layer. Any errors from the database
are automatically caught by a global exception filter and translated into
standard, user-friendly HTTP error responses. This uniform handling prevents
inconsistencies and keeps each module simple and focused.

Table~\ref{tab:modules} lists the feature modules completed so far.

\begin{table}[ht]
	\centering
	\caption{Implemented backend feature modules.}
	\label{tab:modules}
	\resizebox{\textwidth}{!}{%
		\begin{tabular}{lll}
			\hline
			\textbf{Module} & \textbf{Status} & \textbf{Responsibilities / key endpoints} \\ \hline
			\texttt{auth} & \cmark & Register, login, refresh, logout, \texttt{GET /auth/me}; JWT + refresh rotation \\
			\texttt{user} & \cmark & Identity service (find/create user, verify password); used by \texttt{auth} \\
			\texttt{provider-profile} & \cmark & Profile CRUD and skill attachment under \texttt{/provider-profiles} \\
			\texttt{skills} & \cmark & Public catalogue: \texttt{GET /skills/categories}, \texttt{GET /skills} \\
			\texttt{health} & \cmark & Liveness / readiness probe (\texttt{GET /health}) \\
			\texttt{metric} & \cmark & Metrics endpoint \\
			\texttt{service-request} & \xmark & Phase~3, not started \\
			\texttt{matching} & \xmark & Phase~4, not started \\
			\texttt{engagement / rating} & \xmark & Phase~5, not started \\
			\hline
		\end{tabular}%
	}
\end{table}

\subsection{Frontend Module Pattern}
The frontend mirrors the backend's modular organization. Each feature module owns
its API calls, TanStack Query cache-key factories (functions that generate cache
identifiers), custom hooks (split into query hooks for fetching and mutation
hooks for creating/updating), feature-specific components, and page definitions.

A deliberate separation distinguishes \textbf{server data} (information from the
API via TanStack Query, treated as a read-through cache) from \textbf{client
state} (managed by Jotai, temporary UI information like whether a menu is open or
a form is being edited). Server data is never copied into client state; instead,
client state only holds ephemeral information. This architectural separation
prevents the common problem where the UI disagrees with what the API says the
data is.

The application uses file-based routing via TanStack Router. Pages are defined as
files in the \texttt{src/routes/} directory, and the framework automatically
generates the navigation structure. Pathless layouts (\texttt{\_authenticated},
\texttt{\_anonymous}) act as route guards, enforcing authentication at the layout
level, so unauthenticated users cannot access protected pages. Role-based gating
sits on individual leaf routes where needed (e.g., only service providers can
access the profile setup page).

Currently, authentication and provider profiles
are fully implemented with their complete user flows. Routes for future phases
(provider discovery, engagements, admin panel) exist as placeholder.

\section{Authorization Model}
\label{sec:rbac-impl}
Authorization follows the role-based model from Chapter~\ref{ch:work}, refined
into \textbf{three layers}:

\begin{enumerate}
	\item \textbf{Role check.} A coarse, route-level guard keyed on the user's
	      role (\texttt{client}, \texttt{service\_provider}, \texttt{admin}). An
	      administrator is allowed implicitly.
	\item \textbf{Permission check.} A fine-grained guard that requires a named
	      permission in canonical \texttt{resource:verb:scope} form (e.g.,
	      \texttt{provider\_profile:read:any} means "read any provider profile").
	      Permissions are the single source of truth in the shared package: one
	      canonical list generates both the Zod validation schema and the named
	      constants used throughout the codebase, and a role-to-permissions map
	      records each role's grants.
	\item \textbf{Ownership check.} A service-layer check that the acting user
	      owns the resource (or is an administrator). This blocks horizontal
	      access to another user's data.
\end{enumerate}

These three layers work together to gate access. The permission catalogue, the
middle layer, is the most granular control, specifying exactly which actions each
role can perform. For example, a client can create and read their own service
requests, but cannot read another client's requests. The system currently defines
thirteen distinct permissions covering provider profiles, service requests,
engagements, ratings, and user management. Table~\ref{tab:rbac} shows which
permissions each role is granted. The administrator role, by default, has access
to everything.

Both the role and permission checks are fully implemented and tested. A
remaining task is to apply the authentication requirement globally so that all
endpoints are protected by default.

\begin{table}[ht]
	\centering
	\caption{Role-to-permission mapping (administrator is implicit).}
	\label{tab:rbac}
	\begin{tabular}{ll}
		\hline
		\textbf{Role} & \textbf{Permissions} \\ \hline
		\texttt{client} &
			\begin{tabular}[t]{@{}l@{}}
				\texttt{provider\_profile:read:any} \\
				\texttt{service\_request:create|read|update|cancel:own} \\
				\texttt{engagement:read:own}, \texttt{rating:create:own}
			\end{tabular} \\ \hline
		\texttt{service\_provider} &
			\begin{tabular}[t]{@{}l@{}}
				\texttt{provider\_profile:create|update:own} \\
				\texttt{engagement:read|update:own}
			\end{tabular} \\ \hline
	\end{tabular}
\end{table}

\section{Data Model}
The built schema centres on identity and provider data. The tables for requests,
matching, engagements, and ratings are designed but not yet created.
Figure~\ref{fig:er} shows the entities that exist today and how they relate.

\begin{figure}[ht]
	\centering
	\begin{tikzpicture}[
		ent/.style={draw, rounded corners, align=left, font=\scriptsize,
			text width=3.1cm, inner sep=4pt, fill=blue!4},
		rel/.style={-, thick},
		card/.style={font=\tiny, fill=white, inner sep=1pt}
	]
		\node[ent] (users) {\textbf{users}\\\hrulefill\\ id (uuid)\\ email, name\\ role, status\\ password\_hash};
		\node[ent] (rt) [right=2.0cm of users] {\textbf{refresh\_tokens}\\\hrulefill\\ id, user\_id\\ family\_id\\ token\_hash\\ expires/revoked};
		\node[ent] (pp) [below=1.4cm of users] {\textbf{provider\_profiles}\\\hrulefill\\ id, user\_id\\ bio, lat, lng\\ rate\_min/max\\ is\_available\\ profile\_status\\ rating\_*};
		\node[ent] (ps) [right=2.0cm of pp] {\textbf{provider\_skills}\\\hrulefill\\ profile\_id\\ skill\_id\\ level\\ years\_of\_exp};
		\node[ent] (sk) [right=2.0cm of ps] {\textbf{skills}\\\hrulefill\\ id\\ name\\ category\_id};
		\node[ent] (sc) [above=1.4cm of sk] {\textbf{skill\_categories}\\\hrulefill\\ id\\ name};

		\draw[rel] (users.east) -- node[card, near start]{1} node[card, near end]{N} (rt.west);
		\draw[rel] (users.south) -- node[card, near start]{1} node[card, near end]{1} (pp.north);
		\draw[rel] (pp.east) -- node[card, near start]{1} node[card, near end]{N} (ps.west);
		\draw[rel] (sk.west) -- node[card, near start]{1} node[card, near end]{N} (ps.east);
		\draw[rel] (sc.south) -- node[card, near start]{1} node[card, near end]{N} (sk.north);
	\end{tikzpicture}
	\caption{Entity-relationship diagram of the currently implemented schema.}
	\label{fig:er}
\end{figure}

Several design decisions ensure data integrity and performance. When users are
deleted, their record is marked as deleted (rather than physically removed) so
that historical references remain valid. A user's identity is kept separate from
the role they hold, so the platform can later support a user with more than one
role without reworking the schema. Provider profiles progress through
explicit states (\textit{draft}, \textit{active}, or \textit{suspended}), which
makes the business logic clear. To support fast matching queries, the database
is indexed to quickly find active, available providers. Geographic location is
stored as latitude and longitude coordinates, and distance calculations are
performed directly in the database using the haversine formula
(Equation~\ref{eq:haversine}) to avoid fetching unnecessary data into memory.

\subsection{Migrations}
The database schema evolves through three versioned migrations, listed in
Table~\ref{tab:migrations}. Each migration is reviewed and applied in order,
ensuring the database structure can be reliably reproduced across deployments.

\begin{table}[ht]
	\centering
	\caption{Database migrations applied to date.}
	\label{tab:migrations}
	\resizebox{\textwidth}{!}{%
		\begin{tabular}{lll}
			\hline
			\textbf{\#} & \textbf{Migration} & \textbf{Effect} \\ \hline
			1 & \texttt{Init} & Create \texttt{users} and \texttt{refresh\_tokens} (rotation chain) \\
			2 & \texttt{AddUserFieldsAndProviderProfileTables} & Add user status / audit fields; create skills, categories, profiles, provider\_skills \\
			3 & \texttt{AddProviderProfileStatus} & Add \texttt{profile\_status}; partial index on active profiles \\
			\hline
		\end{tabular}%
	}
\end{table}

\section{Quality Assurance}
Quality is enforced by an automated testing strategy and a
continuous-integration pipeline.

\subsection{Testing Strategy}
The project uses a \textbf{multi-layered testing approach}. \textbf{Unit tests}
validate individual pieces of code in isolation, with dependencies mocked out to
ensure tests run quickly. \textbf{Integration tests} run against a real database
with realistic data, verifying that all the pieces work together correctly.
\textbf{Coverage thresholds} are enforced to ensure that at least 80\% of each
module's code is exercised by tests, a safeguard against untested edge cases.

Both backend and frontend are tested. Frontend tests also render components on
screen to verify the user interface behaves correctly.

A k6-based load-testing harness also exists and is run manually during
development. Wiring it into the pipeline and establishing baseline measurements
for key flows (such as login, profile retrieval, and, once it is built,
request-to-match) is planned for the evaluation phase.

\subsection{Continuous Integration}
Six automated checks scrutinize every code change before it can be merged to the
main branch, listed in Table~\ref{tab:ci}. A smart orchestrator determines which
checks to run based on what files changed, so developers don't wait for unrelated
tests. For example, if only frontend files changed, the backend integration tests
are skipped.

\textbf{No code can be merged if any check fails.} This ensures that the main
branch remains in a releasable state: code that breaks the build, violates style
guidelines, has security vulnerabilities, or fails tests is caught before it
reaches production.

\begin{table}[ht]
	\centering
	\caption{Continuous-integration workflows.}
	\label{tab:ci}
	\resizebox{\textwidth}{!}{%
		\begin{tabular}{ll}
			\hline
			\textbf{Workflow} & \textbf{What it checks} \\ \hline
			Style & Code formatting and commit message format \\
			Type checking & TypeScript compilation errors \\
			Unit tests & Fast, isolated tests of individual functions \\
			Security audit & Dependencies for known vulnerabilities \\
			Integration tests & Full system behavior with a real database \\
			Deployment readiness & Full system smoke test and startup health \\
			\hline
		\end{tabular}%
	}
\end{table}

The same containerization underpins the live deployment described in
Section~\ref{sec:deploy}. The images that the deployment-readiness check builds
and smoke-tests are the images deployed to the hosting provider, so the gap
between a passing pipeline and a running deployment stays small, and deploying
with Docker proved straightforward.

\subsection{Security Measures}
The security architecture rests on several foundations. \textbf{Token
management:} access tokens are signed with a cryptographic key (so they cannot
be forged) and short-lived (expiring after a few minutes). Refresh tokens are
hashed at rest (not stored in plain text), sent as \texttt{httpOnly} cookies
(inaccessible to JavaScript, preventing XSS theft), and rotated as a family
chain, so if a refresh token is compromised, revoking it revokes the entire chain of
tokens created from it, limiting the window of exposure. \textbf{Passwords:} are
hashed with bcrypt (a one-way function with built-in rate limiting) so that even
if the database is compromised, passwords cannot be recovered. \textbf{Access
control:} the three-layer authorization model described above is enforced
globally.

During the evaluation phase, a comprehensive penetration-testing plan will
validate the security posture across three attack surfaces: authentication
(token forgery, signature stripping, refresh replay attacks, secret strength,
expiry enforcement), authorization (vertical and horizontal privilege
escalation, claim tampering, weight injection), and network concerns (rate
limiting, information disclosure, CORS configuration, and cookie flags). The
plan follows the OWASP Testing Guide and uses the OWASP Top 10 as its reference
list of attack categories. It combines automated scanning (for example OWASP
ZAP, Nikto, and sqlmap) with manual exploration in a controlled staging
environment, and each finding is logged with a severity rating, a reproduction
path, and a remediation that is then re-tested.

These mechanisms, JWT signing and verification, refresh-token rotation, and the
role, permission, and ownership checks, are implemented directly rather than
delegated to an identity framework such as Keycloak, Auth0, or Better Auth. This
is a deliberate decision aligned with the network-engineering focus of the
thesis. Authentication, session lifecycle, and access control are treated as
objects of study, built and reasoned about at the protocol level so they can be
measured and adversarially tested, as the penetration-testing plan above sets
out. A production deployment could adopt a managed identity provider; building
these mechanisms here is part of the contribution.
